\pagetitle{A.2 / REPRESENTATION}{Một mã, hai loại thông tin}{R1 mở rộng: phép đo vùng + khoảng cách tới biên}
\diagram{A2-representation}{Vùng và biên là hai kiểu phép đo trên cùng tọa độ giải phẫu. Teacher gợi ý ID khi học; chỉ $\hat q_{S,j}$ được ghi vào mã. Hình không biểu diễn fusion trực tiếp các tensor teacher/student.}
\point{Địa chỉ chung.} Mỗi ID xác định kiểu đo, lớp mô $c$, vị trí/vùng và tỉ lệ; không chứa tên encoder. Cùng ID luôn có cùng nghĩa dù feature đến từ nhánh nào:
\[
j=(\text{kiểu},c,\text{vị trí/vùng},\text{tỉ lệ}),\qquad
R=(J,\{\hat q_{S,j}\}_{j\in J}).
\]
\point{Hai ví dụ cụ thể.} Với vùng $V_j$, điểm $p_j$ và lớp $c$, target từ nhãn TRAIN là:
\[
q^{R*}_j=\frac{1}{|V_j|}\int_{V_j}\!\mathbf 1[y(p)=c],dp,
\qquad q^{B*}_j=\operatorname{clip}\!\left(\frac{d_c(y,p_j)}{r},-1,1\right).
\]
$q^{R*}$ là tỉ lệ mô; $d_c$ là khoảng cách có dấu tới biên, âm bên trong và dương bên ngoài. Tọa độ/khoảng cách dùng đơn vị vật lý; $r>0$ cố định trước đánh giá. Lớp vắng: $q^{R*}=0$, quy ước $q^{B*}=+1$.
\point{Từ feature sang mã.} Head $H_S(F_S,j)$ dự đoán giá trị student theo ID. Teacher vùng lấy phép đo từ xác suất lớp; teacher biên lấy từ trường khoảng cách. Giữ riêng kiểu đo và chuẩn hóa lỗi theo kiểu, không ép hai giá trị khác nghĩa bằng nhau. Bank vẫn có thể chứa tương phản và moment.
\point{Decoder đọc gì?} Chỉ ID, giá trị và tọa độ truy vấn $p$; không có feature đi vòng. Tính cả ID và độ chính xác giá trị vào dung lượng. Điểm nghiên cứu là mã thích nghi giữ được thông tin vùng/biên hữu ích đến đâu; chưa có bảo đảm không mất thông tin.
